\chapter{复验指南与附加数据解释}
\label{app:reproduction}

\section{环境准备}

以当前仓库根目录为工作目录，准备 RV64 交叉工具链、QEMU、Python 3 与项目 Makefile 需要的基础命令。首先记录代码 commit、工具链版本、QEMU 版本、Host OS/CPU 和命令行参数。如果使用 Windows，建议在 WSL 中运行 make/QEMU，并保证同一 worktree 同时只有一个写共享 image 的构建。

仓库外 API key 只供 live provider 入口使用。\code{agentos-nexus-replay} 不读 key，也不联网。复验时不应将 key 打印到命令、串口日志或报告中。

\section{分层构建}

\begin{lstlisting}[language=bash,caption={建议的静态、编译与 Guest 回归顺序}]
# 1. Fast local contracts
make local-host-selftests
make agent-uapi-check TOOLPREFIX=riscv64-linux-gnu-

# 2. Production module and stack gates
make agent-module-check TOOLPREFIX=riscv64-linux-gnu-
make build TOOLPREFIX=riscv64-linux-gnu-
make kernel-stack-check TOOLPREFIX=riscv64-linux-gnu-
make user-stack-check TOOLPREFIX=riscv64-linux-gnu-

# 3. Guest functional suite
make agentos-test TOOLPREFIX=riscv64-linux-gnu-
\end{lstlisting}

快速 checker 失败时先修复合同，不要直接启动 QEMU 观察“是否刚好通过”。构建通过后仍不能声称功能通过；必须运行对应 Guest 场景并检查负路径 marker。

\section{Nexus 产品复验}

\subsection{静态合同}

\begin{lstlisting}[language=bash]
make agentos-nexus-check
make agentos-nexus-image TOOLPREFIX=riscv64-linux-gnu-
\end{lstlisting}

第一条命令检查 Host profile/schema、N1 Guest loop、artifact/provenance、审批和 observer 来源边界，不启动 QEMU。第二条命令构建 RV64 Guest 与镜像，并在 mkfs 前检查 Nexus raw binary 不超过 274432 bytes。

\subsection{固定 replay 闭环}

\begin{lstlisting}[language=bash]
make agentos-nexus-replay TOOLPREFIX=riscv64-linux-gnu-
\end{lstlisting}

成功不只是 Guest 打印 \code{passed}。应同时保留 controller、observer 和 daemon 输出，并确认：

\begin{itemize}
  \item HELLO 的 \code{profile=nexus} 与 \code{features=task_event_v1}正确；
  \item 四个业务 Agent 都拥有正数且唯一的 control identity；
  \item System 完成无输入快照并发布 artifact；
  \item Research 先有一个 FAILED terminal，后有新 task 成功，失败 task 不被伪改为成功；
  \item Analyst report 同时绑定 System/Research handle 与 payload SHA-256；
  \item \code{publish_report} 的审批决定字段与请求完全对齐，deny 后内核 \code{artifact_update} 零执行；
  \item observer 中 TASK/audit/correlation/Context 可对齐，正常关闭后无新 telemetry。
\end{itemize}

对 digest-bound fixture 进行修改时，必须从当前 Guest 真实 \code{MODEL_REQUEST} 捕获新 digest，不得使用通配符、空 fixture 或 Host 归一化掩盖请求漂移。修改 history/model projection 后，至少连续两个 fresh strict boot 通过，才能说明该 fixture 对当前源码可复现。

\subsection{现场交互}

\begin{lstlisting}[language=bash]
# Window 1: daemon + controller
make agentos-nexus TOOLPREFIX=riscv64-linux-gnu-

# Window 2: read-only observer
make agentos-nexus-observe
\end{lstlisting}

左窗口提交用户目标并处理审批，右窗口展示内核事实。不要同时启动第二个会写共享 image/runtime 的 Nexus daemon。使用 \code{/quit} 正常关闭，不删除 runtime 文件伪装会话结束。

\section{Contest demo 与 paired 证据}

\begin{lstlisting}[language=bash]
make contest-demo-check
make contest-demo TOOLPREFIX=riscv64-linux-gnu-
\end{lstlisting}

采集时必须保留同一 AgentOS Guest 的 traversal/indexed 输出、比较摘要、负载计数、target/order/seed 和 Host 环境。单次 AB 或 BA 不能提供顺序平衡；对外数字应以 paired boot 为单位。Plain uCore 与 AgentOS-uCore 的双目标比较由 \code{dual-platform-run} 负责，交互式 Nexus 也不能代替这个确定性基准。

\section{一次性数据集的离线检查}

\begin{lstlisting}[language=bash]
python one_shot_metrics/validate.py \
  --tables one_shot_metrics/data/20260811/tables \
  --output ../agentos-20260811-reproduced/validation.json

python one_shot_metrics/plot.py \
  --tables one_shot_metrics/data/20260811/tables \
  --output-dir ../agentos-20260811-reproduced/figures \
  --format png,pdf
\end{lstlisting}

两条命令只做离线验证和重绘，不启动 QEMU。输出应放到仓库外 scratch，不覆盖 \code{one_shot_metrics/data/20260811}。规范目录包含 \code{COMPLETED} 和 145-entry manifest inventory，日常开发不会重新采集它。

\section{附加 I/O 混合 scope 图}

\ReportFigure[0.96\textwidth]
  {../figures/performance/06_normalized_io_heatmaps.pdf}
  {../figures/performance/06_normalized_io_heatmaps.pdf}
  {内核 I/O 和工作量归一化热力图；owner/window/scope 不同，仅作附加工程视图}
  {fig:io-heatmap-appendix}

该图刻意放在附录。\code{bytes_read} 是 workflow lane/core 报告的 total，其他 contest numerator 主要是 shared global/end-to-end delta，Task 则是直接 work count/bytes。它们的 owner、window 和 denominator 不同，不能混成单一进程 I/O 归因。热力图用于快速查找工作量差异，不是主性能结论。

\section{核心数字复核表}

\begin{table}[htbp]
  \centering
  \caption{正文引用的核心数字及其正确解释}
  \label{tab:canonical-numbers}
  \begin{tabularx}{\textwidth}{L{4.0cm}L{3.8cm}Y}
    \toprule
    项目 & 规范值 & 解释约束 \\
    \midrule
    配对 workflow & 16 fresh QEMU boot & boot 是独立单位，AB/BA 各占一半 \\
    core 中位时间 & 34.713/13.294 ms & traversal/indexed core，不是纯 syscall latency \\
    core 配对加速 & 3.118$\times$，16/16 wins & 仅对当前 core 窗口 \\
    E2E 配对 & 3/16 wins，$+13.452$ ms & positive 表示 indexed 窗口更慢 \\
    AB/BA 顺序分层 & 1.4535$\times$/5.4805$\times$ & 提示强顺序/boot 层影响 \\
    catalog grid & $3\times4\times15=180$ pair & 15 是同 boot inner pair，不是 15 boot \\
    Task sequences & 96，4 boot$\times$8 round & 每序列 16 op \\
    Task 中位数 & 561/2051/1620.5 us & compact batch/scalar V3/SQ-CQ 完整 sequence \\
    EEVDF wake & 504 probe；425 zero、79 one tick & exact trace，不从直方图插值 \\
    Jain 中位数 & 1.000000，0.999985，0.999993，0.999985 & 并发 1/2/3/4；1 是平凡基线 \\
    数据集规模 & 30 boot, 33 raw, 19 tables, 7,498 rows & 一次性工程采集，不是跨硬件总体推断 \\
    \bottomrule
  \end{tabularx}
\end{table}

\section{报告工程的重建}

\begin{lstlisting}[language=bash]
cd docs/final-report
latexmk -xelatex -interaction=nonstopmode -halt-on-error main.tex
cp main.pdf AgentOS-uCore-final-report.pdf
\end{lstlisting}

若环境中的 \code{latexmk} 不可用，依次执行一次 XeLaTeX、一次 BibTeX 和两次
XeLaTeX，再复制最终命名 PDF，生成结果等价。

构建后应检查：所有图文件存在；无 undefined reference/citation；无 overfull box；所有页面在 A4 范围内；图题、表题、页眉和页码不重叠；连续页面无空白占位图。最终 PDF 固定输出为 \code{docs/final-report/AgentOS-uCore-final-report.pdf}。
